\documentclass[12pt]{article}
\usepackage[utf8]{inputenc}
\usepackage[T1]{fontenc}
\usepackage{hyperref}
\title{Ruma Language Resources for K-12 Curriculum Development}
\author{Compiled by Claude Sonnet 5.5 (claude-sonnet-5-5) for Andr\'as Kornai}
\date{}
\begin{document}
\maketitle

\section{Introduction}
Ruma (ISO 639-3 \texttt{ruz}) is a Kainji language of Kaduna State, Nigeria. It is also called Ruruma or Turama. Wikipedia gives about 5,100 speakers (a 2000 figure), and Joshua Project gives about 17,000. Joshua Project also notes heavy Hausa influence, with many words replaced by Hausa ones. Some Endangered Languages Catalogue entries titled ``Ruma'' refer to the Turama River in Papua New Guinea and probably describe a different language, so they are not used here.

\section{Primary Educational Resources}
No primers, textbooks or school programmes were found.

\section{Linguistic and Reference Materials}
\begin{itemize}
\item Wikipedia, ``Ruma language'': a short overview with classification and speaker numbers. \url{https://en.wikipedia.org/wiki/Ruma_language}
\end{itemize}
No dictionary, grammar or orthography description was found.

\section{Religious and Community Resources}
\begin{itemize}
\item Joshua Project, ``Ruruma in Nigeria'': people-group profile with alternate names, a speaker estimate and community history. \url{https://joshuaproject.net/people_groups/14596}
\item Joshua Project language page for \texttt{ruz}. \url{https://joshuaproject.net/index.php/languages/ruz}
\end{itemize}

\section{Digital and Audio Resources}
No audio, corpus or software resources were found.

\section{Curriculum Development Context}
The language is not known to be taught in schools. Very little documentation is available online, and Hausa influence is strong, so the language may be shifting toward Hausa. A native-language curriculum would first need basic documentation work and community consultation.
\end{document}